% Generated from locale/tr/content/set-theory/spine/rank.tex; do not hand-edit.


\olfileid[tr]{sth}{spine}{rank}
\olsection{Rank}

Aşamaları $V_\alpha$'lar olarak tanımladığımıza ve her kümenin bir aşamanın
alt kümesi olduğunu bildiğimize göre bir kümenin \emph{rankını}
tanımlayabiliriz. Sezgisel olarak $A$'nın rankı, $A$'nın oluşturulduğu ilk
andır. Daha kesin olarak:

\begin{defn}\ollabel{defnsetrank}
Her $A$ kümesi için $\setrank{A}$, $A\subseteq V_\alpha$ koşulunu sağlayan
en küçük $\alpha$ sıra sayısıdır.
\end{defn}
\begin{prop}\ollabel{ranksexist}
	Her $A$ için $\setrank{A}$ vardır.
\end{prop}
\begin{proof}
	Alıştırma olarak bırakılmıştır.
\end{proof}
\begin{prob}
	\olref[sth][spine][rank]{ranksexist} sonucunu ispatlayın.
\end{prob}
\noindent
Rankların iyi sıralanması bazı önemli sonuçları ispatlamamızı sağlar:
\begin{prop}\ollabel{valphalowerrank}
Her $\alpha$ sıra sayısı için
$V_\alpha=\Setabs{x}{\setrank{x}\in\alpha}$.
\end{prop}

\begin{proof}
$\setrank{x}\in\alpha$ ise
$x\subseteq V_{\setrank{x}}\in V_\alpha$ olur; $V_\alpha$ kuvvetli
olduğundan (\olref[Valphabasic]{Valphabasicprops} birkaç kez kullanılarak)
$x\in V_\alpha$ elde edilir. Tersine $x\in V_\alpha$ olsun. Basit sonlu
ötesi tümevarımla şu görülür: $\alpha$ ardıl $\ordsucc{\beta}$ ise
$x\subseteq V_\beta$ ve dolayısıyla $\setrank{x}\leq\beta<\alpha$ olur;
$\alpha$ limit ise bir $\beta<\alpha$ için $x\in V_\beta$ ve dolayısıyla
yine $\setrank{x}<\alpha$ olur. Böylece $\setrank{x}\in\alpha$.
\end{proof}
\begin{prob}
	\olref[sth][spine][rank]{valphalowerrank} içinde anılan basit sonlu ötesi
	tümevarımı tamamlayın.
\end{prob}

\begin{prop}\ollabel{rankmemberslower}
$B\in A$ ise $\setrank{B}\in\setrank{A}$.
\end{prop}

\begin{proof}
\olref{valphalowerrank} gereği
$A\subseteq V_{\setrank{A}}=
\Setabs{x}{\setrank{x}\in\setrank{A}}$ olur.
\end{proof}
\noindent
Bu olguyla, bir tür tümevarım kullanarak \emph{bütün kümeler} hakkında
sonuç ispatlamamızı sağlayan bir sonuç kurabiliriz:

\begin{thm}[$\in$-Tümevarım Şeması]
Her $\phi$ formülü için:
\[
	\forall A((\forall x \in A)\phi(x) \lif \phi(A)) \lif \forall A \phi(A).
\]
\end{thm}

\begin{proof}
Karşıt-tersi ispatlayacağız. $\lnot\forall A\phi(A)$ olduğunu varsayın.
Sonlu Ötesi Tümevarım
(\olref[ordinals][basic]{ordinductionschema}) gereği olabilecek en küçük
ranka sahip bir $\phi$-olmayan vardır; yani $\lnot\phi(A)$ ve
$\forall x(\setrank{x}\in\setrank{A}\lif\phi(x))$ olan bir $A$ vardır.
Şimdi $x\in A$ ise \olref{rankmemberslower} gereği
$\setrank{x}\in\setrank{A}$ ve dolayısıyla $\phi(x)$ olur. Böylece
$(\forall x\in A)\phi(x)\land\lnot\phi(A)$ elde edilir.
\end{proof}\noindent
Bu güçlü sonucu gayriresmî olarak şöyle açıklayabiliriz. Bir kümede bulunan
her !!{element} $\phi$ olduğunda kümenin kendisi de $\phi$ ise $\phi$'ye
\emph{kalıtsal} diyelim. O zaman $\in$-Tümevarım şunu söyler: $\phi$
kalıtsalsa her küme $\phi$'dir.

Rank tartışmasını (şimdilik) tamamlamak için daha önce birkaç kez haber
verdiğimiz bazı iddiaları ispatlayacağız.

\begin{prop}\ollabel{ranksupstrict}
$\setrank{A} = \supstrict_{x \in A}\setrank{x}$.
\end{prop}

\begin{proof}
$\alpha=\supstrict_{x\in A}\setrank{x}$ olsun.
\olref{rankmemberslower} gereği $\alpha\leq\setrank{A}$ olur. Fakat
$x\in A$ ise $\setrank{x}\in\alpha$ ve dolayısıyla
\olref{valphalowerrank} gereği $x\in V_\alpha$ olur. Böylece
$A\subseteq V_\alpha$, yani $\setrank{A}\leq\alpha$ olur. Sonuç olarak
$\setrank{A}=\alpha$.
\end{proof}

\begin{cor}\ollabel{ordsetrankalpha}
Her $\alpha$ sıra sayısı için $\setrank{\alpha}=\alpha$.
\end{cor}

\begin{proof}
Sonlu ötesi tümevarım için bütün $\beta\in\alpha$ üzerinde
$\setrank{\beta}=\beta$ olduğunu varsayın. Şimdi
$\setrank{\alpha}=\supstrict_{\beta\in\alpha}\setrank{\beta}=
\supstrict_{\beta\in\alpha}\beta=\alpha$ olur; burada
\olref{ranksupstrict} kullanılmıştır.
%	First note that $\setrank{\alpha} \neq \beta$ for any $\beta \in
%	\alpha$. For otherwise we would have $\beta = \setrank{\beta} \in
%	\setrank{\alpha} = \beta$ by \olref{rankmemberslower}, i.e.,
%	$\beta \in \beta$, a contradiction. 
%
%	Now note that $\alpha \subseteq V_\alpha$. For $\alpha =
%	\Setabs{\beta}{\beta \in \alpha}$ by
%	\olref[sfr][ordinals][basic]{ordissetofsmallerord}. And if
%	$\beta \in \alpha$ then $\beta \subseteq V_{\setrank{\beta}} =
%	V_\beta \in V_\alpha$, hence $\beta \in V_\alpha$, by
%	\olref[sfr][spine][Valphabasic]{Valphabasicprops} twice. 
\end{proof}

Son olarak \olref[foundation]{sec} sonunda vaat edilen, $\ZFminus$'ın
$\emph{Düzenlilik}\Rightarrow\emph{Temellendirme}$ koşullusunu
ispatladığı sonucunun hızlı bir ispatını verelim. (``Rank'' kavramı ile
\olref{rankmemberslower} bu ispatta kullanılabilir; çünkü bu kesimin
başında belirtildiği gibi bunlar $\ZFminus+\text{Düzenlilik}$ kullanılarak
sunulabilir.)

\begin{prop}[$\ZFminus + \text{Düzenlilik}$ içinde çalışarak]\ollabel{zfminusregularityfoundation}
Temellendirme geçerlidir.
\end{prop}

\begin{proof}
$A\neq\emptyset$ sabit olsun ve $A$ içinde olabilecek en küçük ranka sahip
bir $B$ seçin. $c\in B$ ise \olref{rankmemberslower} gereği
$\setrank{c}\in\setrank{B}$ olur; dolayısıyla $B$'nin seçimi gereği
$c\notin A$'dır. Böylece $A\cap B=\emptyset$ olur.
\end{proof}

